% =====================================================================
\chapter{Миксеры}
\label{ch:mixes}
% =====================================================================
\chapterintro
  {как входы попадают в каналы, что делают вес, смещение, режимы сложения, задержка и
   замедление, как собрать элевоны, V-хвост, флапероны и дифференциал элеронов}
  {модель с настроенными входами}

\section{Страница Mixes}

\mpath{MDL\sep Mixes} показывает список каналов \sw{CH1}…\sw{CH32} и строки миксеров под ними.
Значение канала --- результат всех его строк. У новой модели миксеры простые: каждый вход
идёт в свой канал с весом 100\,\%.

\begin{center}
\lcd{\inv{MIXES}\hfill 4/64\\
CH1\quad {[I1]}Ail\quad 100\\
CH2\quad {[I2]}Ele\quad 100\\
CH3\quad {[I3]}Thr\quad 100\\
CH4\quad {[I4]}Rud\quad 100\\
CH5}
\end{center}

Чтобы добавить строку, встаньте на канал и нажмите \btn{ENTER} (для пустого канала) или
долгое \btn{ENTER} → \menu{Insert after}.

\section{Параметры строки миксера}

\begin{center}\small
\renewcommand{\arraystretch}{1.2}
\begin{tabularx}{\linewidth}{@{}p{30mm}X@{}}
\toprule
\menu{Mix name} & имя строки \\
\menu{Source} & источник: вход (\sw{I1}…), стик напрямую, переключатель, \sw{MAX},
  другой канал, \sw{GV}, телеметрия \\
\menu{Weight} & вес $-500\ldots+500\,\%$ (отрицательный --- инверсия) \\
\menu{Offset} & смещение \\
\menu{Trim} & учитывать триммер источника \\
\menu{Curve} & кривая или функция (как во входах) \\
\menu{Flight modes} & в каких полётных режимах активна строка \\
\menu{Switch} & условие активности \\
\menu{Warning} & звуковое предупреждение, пока строка активна \\
\menu{Multiplex} & как объединять со строками выше: \menu{Add} --- сложить,
  \menu{Multiply} --- умножить, \menu{Replace} --- заменить \\
\menu{Delay up/down} & задержка реакции на изменение \\
\menu{Slow up/down} & скорость изменения значения (плавность), в секундах на полный ход \\
\bottomrule
\end{tabularx}
\end{center}

\subsection{Как считается канал}

Строки канала обрабатываются \emph{сверху вниз}. Каждая активная строка вычисляет своё
значение $v_i = \text{curve}(\text{source}) \cdot \text{weight} + \text{offset}$ и объединяет
его с накопленным результатом:
\[
  \text{Add: } r \gets r + v_i, \qquad
  \text{Multiply: } r \gets r \cdot v_i, \qquad
  \text{Replace: } r \gets v_i.
\]
Неактивные строки (переключатель не включён, другой полётный режим) пропускаются.

\section{Вспомогательные каналы (AUX)}

Самое частое применение --- переключатели на каналы 5 и выше:

\begin{center}\small
\begin{tabular}{@{}llll@{}}
\toprule
Канал & Source & Weight & Назначение \\
\midrule
\sw{CH5} & \sw{SA} & 100 & арм (для ELRS --- обязательно двухпозиционный) \\
\sw{CH6} & \sw{SB} & 100 & полётные режимы контроллера \\
\sw{CH7} & \sw{SC} & 100 & бипер, турель, запись \\
\sw{CH8} & \sw{S1} & 100 & наклон камеры \\
\bottomrule
\end{tabular}
\end{center}

Двухпозиционный переключатель как источник даёт $-100$ или $+100$, а трёхпозиционный ---
$-100$/$0$/$+100$.

\section{Элевоны (летающее крыло)}

У летающего крыла нет отдельного руля высоты: две поверхности-элевона работают и как
элероны (в разные стороны), и как руль высоты (в одну сторону).

\begin{figure}[H]
\centering
\begin{tikzpicture}[node distance=4mm and 12mm]
  \node[blk,fill=blkin] (ail) {\sw{I1} Ail};
  \node[blk,fill=blkin,below=10mm of ail] (ele) {\sw{I2} Ele};
  \node[blk,fill=blkout,right=35mm of ail] (ch1) {\sw{CH1} левый элевон};
  \node[blk,fill=blkout,right=35mm of ele] (ch2) {\sw{CH2} правый элевон};
  \draw[arr] (ail.east) -- node[above,lbl,pos=0.5]{$+50\,\%$} (ch1.west);
  \draw[arr] (ele.east) -- node[below,lbl,pos=0.12,sloped]{$+50\,\%$} (ch1.west);
  \draw[arr] (ail.east) -- node[above,lbl,pos=0.12,sloped]{$-50\,\%$} (ch2.west);
  \draw[arr] (ele.east) -- node[below,lbl,pos=0.5]{$+50\,\%$} (ch2.west);
\end{tikzpicture}
\caption{Миксер элевонов}
\end{figure}

\begin{recipe}{элевоны}
\begin{center}\small
\begin{tabular}{@{}llll@{}}
\toprule
Канал & Source & Weight & Multiplex \\
\midrule
\sw{CH1} & \sw{I1} Ail & $+50$ & Add \\
         & \sw{I2} Ele & $+50$ & Add \\
\sw{CH2} & \sw{I1} Ail & $-50$ & Add \\
         & \sw{I2} Ele & $+50$ & Add \\
\sw{CH3} & \sw{I3} Thr & $100$ & Add \\
\bottomrule
\end{tabular}
\end{center}
Веса по 50\,\% нужны, чтобы при одновременном полном крене и тангаже сумма не превысила
100\,\%. Направления проверяются на земле: ручку на себя --- оба элевона вверх; ручку
вправо --- правый элевон вверх, левый вниз. Если какая-то поверхность ходит неверно,
инвертируйте её в \menu{Outputs} или поменяйте знак веса.
\end{recipe}

\section{V-образное оперение}

Аналогично элевонам, но смешиваются руль высоты и руль направления:
\sw{CH2} = Ele $+50$ и Rud $+50$; \sw{CH4} = Ele $+50$ и Rud $-50$. Ручка на себя ---
обе поверхности вверх; правая педаль --- обе поверхности вправо (если смотреть сзади).

\section{Флапероны и дифференциал элеронов}

Если каждый элерон на своём канале (\sw{CH1} --- левый, \sw{CH6} --- правый), можно:
\begin{itemize}
  \item \textbf{флапероны}: оба элерона вниз как закрылки. Строка
        \sw{CH1}: \sw{SC} вес $-30$, \sw{CH6}: \sw{SC} вес $+30$ (знаки зависят от установки серв);
  \item \textbf{дифференциал}: элерон, идущий вниз, отклоняется меньше, чем идущий вверх ---
        это уменьшает неблагоприятное рыскание. В строке с источником \sw{I1} задайте
        \menu{Curve} = \menu{Diff} 30\,\%.
\end{itemize}

\section{Отсечка газа через Replace}

\begin{recipe}{отсечка газа (throttle cut)}
В канале газа \sw{CH3} после основной строки добавьте:
Source = \sw{MAX}, Weight = $-100$, Switch = \sw{SD↓}, Multiplex = \menu{Replace}.
Пока \sw{SD} внизу, канал газа \emph{заменяется} на $-100\,\%$ независимо от стика.
То же можно сделать специальной функцией \menu{Override} (глава~\ref{ch:special}).
\end{recipe}

\section{Замедление (Slow) и задержка (Delay)}

\menu{Slow up/down} задаёт время прохождения полного хода. Примеры:
\begin{itemize}
  \item шасси и закрылки --- 2--3 секунды, чтобы они выпускались плавно, «как у настоящих»;
  \item плавный старт мотора у планера --- \menu{Slow up} = 1\,с на канале газа.
\end{itemize}
\menu{Delay} задерживает реакцию: например, створки шасси открываются сразу, а сами
стойки выходят через секунду.

\section{Компенсация тангажа при выпуске закрылков}

При выпуске закрылков самолёт обычно задирает нос. Добавьте в канал руля высоты строку:
Source = \sw{SC} (переключатель закрылков), Weight = $-10$ (подбирается в полёте),
\menu{Curve} = \menu{x>0} --- компенсация будет только при выпущенных закрылках.

\begin{summary}
\begin{itemize}
  \item Миксеры связывают входы с каналами; строки канала считаются сверху вниз.
  \item \menu{Add} --- сложить, \menu{Multiply} --- умножить, \menu{Replace} --- заменить.
  \item Элевоны и V-хвост --- два входа с весами по 50\,\% и разными знаками.
  \item \menu{Slow} даёт плавность, \menu{Delay} --- задержку.
\end{itemize}
\end{summary}
